\section{Kỹ thuật đặc trưng (Feature Engineering)}
\label{sec:feature_eng}

\subsection{Mã hóa đặc trưng High Cardinality --- Target Encoding}
\label{subsec:target_encoding}

Ba đặc trưng phân loại \texttt{City}, \texttt{Profession} và \texttt{Degree} có số lượng category rất lớn (hàng chục đến hàng trăm giá trị unique), khiến One-Hot Encoding không khả thi vì sẽ tạo quá nhiều cột thưa. Nhóm tác giả áp dụng \textbf{Target Encoding} --- thay thế mỗi category bằng giá trị trung bình của biến mục tiêu trong category đó, kết hợp smoothing để giảm overfitting:

\begin{equation}
\text{TargetEnc}(c) = \frac{\displaystyle\sum_{i:\, x_i = c} y_i + \lambda \cdot \bar{y}}{n_c + \lambda}
\label{eq:target_encoding}
\end{equation}

trong đó:
\begin{itemize}[noitemsep]
    \item $c$: category cần mã hóa
    \item $n_c$: số lần category $c$ xuất hiện trong tập dữ liệu
    \item $\bar{y}$: giá trị trung bình của biến mục tiêu trên toàn bộ tập dữ liệu (global mean)
    \item $\lambda$: hệ số smoothing, kiểm soát mức độ ``co về'' global mean
\end{itemize}

Khi $n_c$ nhỏ (category hiếm), giá trị encoding bị kéo về $\bar{y}$ nhiều hơn, giúp tránh overfitting. Để ngăn ngừa data leakage, Target Encoding được thực hiện trong vòng lặp Cross-Validation: chỉ fit trên tập train của mỗi fold và transform trên tập validation.

\subsection{Phân nhóm tuổi (Age Binning)}
\label{subsec:age_binning}

Dựa trên phân tích EDA (Mục~\ref{subsec:bivariate}), nhóm tác giả nhận thấy mối quan hệ giữa tuổi và trầm cảm không tuyến tính đơn thuần mà có sự khác biệt rõ rệt theo nhóm tuổi. Do đó, kỹ thuật \textit{Discretization/Binning} được áp dụng để chia \texttt{Age} thành 4 nhóm ordinal:

\begin{itemize}[noitemsep]
    \item Nhóm 0: 18--29 tuổi (tỷ lệ trầm cảm 57,0\%)
    \item Nhóm 1: 30--42 tuổi (tỷ lệ trầm cảm 12,8\%)
    \item Nhóm 2: 43--51 tuổi (tỷ lệ trầm cảm 1,7\%)
    \item Nhóm 3: 52+ tuổi (tỷ lệ trầm cảm 0,3\%)
\end{itemize}

Hình~\ref{fig:age_group} minh họa sự khác biệt rất lớn về tỷ lệ trầm cảm giữa các nhóm tuổi: nhóm 18--29 có tỷ lệ gấp gần 200 lần so với nhóm 52+, đồng thời số lượng mẫu phân bố tương đối đồng đều giữa 4 nhóm (~32.000--38.000 mẫu/nhóm).

\begin{figure}[H]
\centering
\includegraphics[width=0.95\textwidth]{feature_eng/age_group_vs_depression.png}
\caption{Tỷ lệ Depression và số lượng mẫu theo nhóm tuổi (Age Binning)}
\label{fig:age_group}
\end{figure}

\subsection{Tạo đặc trưng tương tác (Interaction Features)}
\label{subsec:interaction_features}

Dựa trên kết quả phân tích đa biến ở Mục~\ref{subsec:multivariate}, nhóm tác giả tạo thêm các đặc trưng kết hợp nhằm nắm bắt hiệu ứng tương tác giữa nhiều yếu tố:

\textbf{Stress Load} --- đo mức tải stress tổng hợp từ áp lực và căng thẳng tài chính:
\begin{equation}
\text{stress\_load} = \text{Pressure} \times \text{Financial\_Stress}
\label{eq:stress_load}
\end{equation}

\textbf{Wellbeing Index} --- chỉ số sức khỏe tổng thể dựa trên mức hài lòng và giấc ngủ:
\begin{equation}
\text{wellbeing\_index} = \frac{\text{Satisfaction} + \text{Sleep\_Duration}}{2}
\label{eq:wellbeing}
\end{equation}

\textbf{Stress-Satisfaction Ratio} --- tỷ lệ giữa áp lực và mức hài lòng:
\begin{equation}
\text{stress\_satisfaction\_ratio} = \frac{\text{Pressure}}{\text{Satisfaction} + 1}
\label{eq:stress_ratio}
\end{equation}

\textbf{Risk Score} --- chỉ số rủi ro cộng gộp từ hai yếu tố cảnh báo mạnh nhất:
\begin{equation}
\text{risk\_score} = \text{Suicidal\_Thoughts} + \text{Family\_History}
\label{eq:risk_score}
\end{equation}

Hình~\ref{fig:interaction_cont} minh họa phân phối các đặc trưng tương tác liên tục và mối quan hệ với Depression. Nhóm trầm cảm có giá trị \texttt{stress\_load} và \texttt{stress\_satisfaction\_ratio} cao hơn rõ rệt, trong khi \texttt{wellbeing\_index} thấp hơn --- phù hợp với kỳ vọng.

\begin{figure}[H]
\centering
\includegraphics[width=0.95\textwidth]{feature_eng/interaction_continuous_vs_depression.png}
\caption{Phân phối và boxplot các đặc trưng tương tác liên tục theo Depression}
\label{fig:interaction_cont}
\end{figure}

Hình~\ref{fig:interaction_disc} cho thấy đặc trưng \texttt{risk\_score} có khả năng phân biệt tốt: khi score = 0, tỷ lệ trầm cảm chỉ 4,7\%; khi score = 1 là 17,9\%; và khi score = 2 (cả hai yếu tố rủi ro đều có) là 32,6\%.

\begin{figure}[H]
\centering
\includegraphics[width=0.95\textwidth]{feature_eng/interaction_discrete_vs_depression.png}
\caption{Tỷ lệ Depression theo poor\_sleep và risk\_score}
\label{fig:interaction_disc}
\end{figure}

\subsection{Lựa chọn đặc trưng (Feature Selection)}
\label{subsec:feature_selection}

Để đánh giá tầm quan trọng của các đặc trưng (bao gồm cả đặc trưng gốc và đặc trưng mới tạo), nhóm tác giả sử dụng \textbf{Random Forest Feature Importance} --- đo mức giảm Gini Impurity trung bình khi split trên mỗi đặc trưng.

Hình~\ref{fig:feat_importance_rf} cho thấy kết quả xếp hạng: \texttt{Age} (0,17) và \texttt{age\_group} (0,13) là hai đặc trưng quan trọng nhất, tiếp theo là \texttt{Suicidal Thoughts} (0,11), \texttt{cgpa\_pressure} (0,10) và \texttt{Profession\_te} (0,07). Đáng chú ý, các đặc trưng tương tác mới tạo (\texttt{stress\_load}, \texttt{cgpa\_pressure}, \texttt{stress\_satisfaction\_ratio}) đều nằm trong top 10, xác nhận giá trị của bước Feature Engineering.

\begin{figure}[H]
\centering
\includegraphics[width=0.7\textwidth]{feature_eng/feat_importance.png}
\caption{Random Forest Feature Importance --- xếp hạng toàn bộ đặc trưng}
\label{fig:feat_importance_rf}
\end{figure}

Kết luận: nhóm tác giả giữ lại toàn bộ đặc trưng có importance $> 0$ (tổng cộng 25 đặc trưng) để đưa vào huấn luyện mô hình, vì các mô hình tree-based có khả năng tự chọn đặc trưng quan trọng trong quá trình học.
